\documentclass[12pt]{article}

\usepackage{fullpage}
\usepackage{mdframed}
\usepackage{multicol}
\usepackage{colonequals}
\usepackage{algpseudocode}
\usepackage{algorithm}
\usepackage[most, breakable]{tcolorbox}
\usepackage[all]{xy}
\usepackage{proof}
\usepackage{mathtools}
\usepackage{bbm}
\usepackage{amssymb}
\usepackage{amsthm}
\usepackage{amsmath}
\usepackage{amsxtra}
\usepackage{tikz-cd}
\usepackage{enumitem}
\usepackage{euler}
\usepackage{mathpazo}
\usepackage{wrapfig}

\newtheorem{theorem}{Theorem}[section]
\newtheorem{theorem*}{Theorem}
\newtheorem{definition}[theorem]{Definition}
\newtheorem{corollary}{Corollary}[theorem]
\newtheorem{lemma}[theorem]{Lemma}
\newtheorem{lemmact}[theorem]{Lemmact}
\newtheorem{proposition}[theorem]{Proposition}
\newtheorem{remark}[theorem]{Remark}
\newtheorem{example}{Example}[section]


% \newtheorem*{exercisehelper}{Exercise.}
% \newenvironment{exercise}[1]{%
%   \IfBlankTF{#1}
%     {\renewcommand{\exercisehelper}{\textbf{Exercise} \unskip}}
%     {\renewcommand\exercisehelper{\textbf{Exercise #1}}}%
%   \exercisehelper
% }{\endexercisehelper}

\theoremstyle{remark}
\newtheorem{exercise}{Exercise}[section]
\newtheorem*{solution}{Solution}
\newtheorem*{notation}{Notation}
\newcommand{\mathcat}[1]{\textup{\textbf{\textsf{#1}}}} % for defined terms

\newenvironment{problem}[1]
{ \begin{tcolorbox}[breakable]\noindent\textbf{Problem #1}.}
{\vskip 6pt \end{tcolorbox}}

\newenvironment{enumalph}
{\begin{enumerate}\renewcommand{\labelenumi}{\textnormal{(\alph{enumi})}}}
{\end{enumerate}}

\newenvironment{enumroman}
{\begin{enumerate}\renewcommand{\labelenumi}{\textnormal{(\roman{enumi})}}}
{\end{enumerate}}

\newcommand{\defi}[1]{\textsf{#1}} % for defined terms



\setlength{\hfuzz}{4pt}

\let\H\relax
\let\P\relax
\newcommand{\bb}{\mathbb}
\newcommand{\A}{\mathbb A}
\newcommand{\H}{\mathbb H}
\newcommand{\E}{\mathbb E}
\newcommand{\P}{\mathbb P}
\newcommand{\C}{\mathbb C}
\newcommand{\I}{\mathbb I}
\newcommand{\K}{\mathbb K}
\newcommand{\N}{\mathbb N}
\newcommand{\Q}{\mathbb Q}
\newcommand{\R}{\mathbb R}
\newcommand{\Z}{\mathbb Z}
\newcommand{\F}{\mathbb F}
\newcommand{\V}{\mathbb V}
\newcommand{\br}{\mathbf{r}}
\newcommand{\RP}{\mathbb{RP}}
\newcommand{\CP}{\mathbb{CP}}
\newcommand{\nbit}[1]{\{0, 1\}^{#1}}
\newcommand{\bits}{\{0, 1\}^{n}}
\newcommand{\bbni}{\bigbreak \noindent}
\newcommand{\norm}[1]{\left\vert\left\vert#1\right\vert\right\vert}
\newcommand{\dbar}{\overline{\partial}}
\let\d\relax
\newcommand{\d}{\partial}
\newcommand{\e}{\epsilon}
\newcommand{\cO}{\mathcal{O}}
\newcommand{\FF}{\mathcal{F}}
\newcommand{\cG}{\mathcal{G}}
\newcommand{\HH}{\mathcal{H}}
\newcommand{\EE}{\mathcal{E}}
\newcommand{\CC}{\mathcal{C}}
\newcommand{\DD}{\mathcal{D}}
\newcommand{\PP}{\mathcal{P}}
\newcommand{\iso}{\cong}
\newcommand{\tism}{Talked to Sair Shaikh'26}
\DeclareMathOperator{\fm}{\mathfrak{m}}
\DeclareMathOperator{\fn}{\mathfrak{n}}
\DeclareMathOperator{\fp}{\mathfrak{p}}
\DeclareMathOperator{\fq}{\mathfrak{q}}
\DeclareMathOperator{\fM}{\mathfrak{M}}
\DeclareMathOperator{\fP}{\mathfrak{P}}


\newcommand{\Frac}{\text{Frac}}
\newcommand{\Tor}{\text{Tor}}
\newcommand{\Ext}{\text{Ext}}
\newcommand{\gr}{\text{gr}}
\newcommand{\pd}{\text{pd}}
\newcommand{\Proj}{\text{Proj}}
\newcommand{\NF}{\text{NF}}
\newcommand{\Jac}{\text{Jac}}
\newcommand{\red}{\text{red}}
\newcommand{\Nil}{\text{Nil}}
\newcommand{\LC}{\text{LC}}
\newcommand{\LM}{\text{LM}}
\newcommand{\LT}{\text{LT}}
\newcommand{\GL}{\text{GL}}
\newcommand{\tp}{\text{T}}
\newcommand{\init}{\text{in}}
\DeclareMathOperator{\Bl}{Bl}


\DeclareMathOperator{\cS}{\mathscr S}
\DeclareMathOperator{\cV}{\mathscr V}
\DeclareMathOperator{\cR}{\mathscr R}
\DeclareMathOperator{\cT}{\mathscr T}
\DeclareMathOperator{\comRing}{\textbf{comRing}}
\DeclareMathOperator{\Ann}{\text{Ann}}
\DeclareMathOperator{\adj}{\text{Adj}}
\DeclareMathOperator{\Top}{\textbf{Top}}
\DeclareMathOperator{\Mod}{\textbf{Mod}}
% \renewcommand{\Top}{\textbf{Top}}
\DeclareMathOperator{\Set}{\textbf{Set}}
\DeclareMathOperator{\alg}{\textbf{alg}}
\DeclareMathOperator{\ord}{\textbf{ord}}
\DeclareMathOperator{\height}{\text{ht}}
\DeclareMathOperator{\trdeg}{\text{tr.deg}}
\DeclareMathOperator{\std}{\text{std}}
\DeclareMathOperator{\depth}{\text{depth}}
\DeclareMathOperator{\ini}{\text{in}}




\let\1\relax
\newcommand{\1}{\mathbf{1}}
\newcommand{\fr}[2]{\left(\frac{#1}{#2}\right)}
\newcommand{\todo}[1]{\textcolor{red}{\textbf{TODO:} #1}}
\newcommand{\vecz}{\mathbf{z}}
\newcommand{\vecr}{\mathbf{r}}
\DeclareMathOperator{\ot}{\leftarrow}
\DeclareMathOperator{\Cinf}{C^{\infty}}
\DeclareMathOperator{\Id}{Id}
\DeclareMathOperator{\Ell}{Ell}
\DeclareMathOperator{\CL}{\mathcal{CL}}
\DeclareMathOperator{\Alt}{Alt}
\DeclareMathOperator{\Aut}{Aut}
\DeclareMathOperator{\ann}{ann}
\DeclareMathOperator{\codim}{codim}
\DeclareMathOperator{\End}{End}
\DeclareMathOperator{\Hom}{Hom}
\DeclareMathOperator{\id}{id}
\DeclareMathOperator{\im}{im}
\DeclareMathOperator{\M}{M}
\DeclareMathOperator{\Mat}{Mat}
\DeclareMathOperator{\Ob}{Ob}
\DeclareMathOperator{\opchar}{char}
\DeclareMathOperator{\opspan}{span}
\DeclareMathOperator{\rk}{rk}
\DeclareMathOperator{\ev}{ev}
\DeclareMathOperator{\sgn}{sgn}
\DeclareMathOperator{\Sym}{Sym}
\DeclareMathOperator{\tr}{tr}
\DeclareMathOperator{\img}{img}
\DeclareMathOperator{\coker}{coker}
\DeclareMathOperator{\Spec}{Spec}
\DeclareMathOperator{\mSpec}{mSpec}
\DeclareMathOperator{\CandE}{CandE}
\DeclareMathOperator{\CandO}{CandO}
\DeclareMathOperator{\argmax}{argmax}
\DeclareMathOperator{\first}{first}
\DeclareMathOperator{\last}{last}
\DeclareMathOperator{\cost}{cost}
\DeclareMathOperator{\dist}{dist}
\DeclareMathOperator{\path}{path}
\DeclareMathOperator{\parent}{parent}
\DeclareMathOperator{\argmin}{argmin}
\DeclareMathOperator{\excess}{excess}
\let\Pr\relax
\DeclareMathOperator{\Pr}{\mathbf{Pr}}
\DeclareMathOperator{\Exp}{\mathbb{E}}
\DeclareMathOperator{\Var}{\mathbf{Var}}
\let\limsup\relax
\DeclareMathOperator{\limsup}{limsup}
%Paired Delims
\DeclarePairedDelimiter\ceil{\lceil}{\rceil}
\let\oldceil\ceil
\renewcommand{\ceil}[1]{\oldceil*{#1}}

\DeclarePairedDelimiter{\floor}{\lfloor}{\rfloor}
\let\oldfloor\floor
\renewcommand{\floor}[1]{\oldfloor*{#1}}


\newcommand{\circled}[2][]{
  \tikz[baseline=(char.base)]{
    \node[shape=circle, draw, inner sep=1.5pt, #1] (char) {$#2$};
  }
}
\newcommand{\dcircled}[2][]{
  \tikz[baseline=(X.base)]{
    \node[circle, draw, double, double distance=2pt, inner sep=1.5pt, #1] (X) {$#2$};
  }
}

\newcommand{\dagstar}{*}

\newcommand{\tbigwedge}{{\textstyle{\bigwedge}}}
\setlength{\parindent}{0pt}
\setlength{\parskip}{5pt}


\usepackage{listings}
\usepackage{courier}
\usepackage{microtype}


\lstset{
  basicstyle=\footnotesize\ttfamily,
  breaklines=true,
  breakatwhitespace=true
  columns=fullflexible,
  keepspaces=true,
  frame=single,
  escapeinside={(*@}{@*)}
}
\usepackage{tikz-cd}
\tikzset{
    lock/.style={circle, draw, minimum size=7mm, inner sep=0pt},
    proc/.style={rectangle, draw, minimum size=6mm, inner sep=0pt},
    unused/.style={draw, dashed, gray},
    hl/.style={line width=1.6pt},
}

\begin{document}

\title{COSC 42: Distributed Computing}

\author{}
\date{}
\maketitle
\vspace{-.8in}


\begin{problem}{1. $n$-Process Locks from Two-Process Locks}

    Reduce $n$-process locks to two process locks. That is, let $n \geq 2$ be any fixed and known integer, and let $L_2$ be a solution to the mutual exclusion problem for two-processes, construct by reduction a lock $L_n$ that solves the mutual exclusion problem for $n$-processes. Your reduction should only use the read and write shared-memory instructions (i.e., it should not use RMW instruction) and it should satisfy the following requirements:
    \begin{itemize}
        \item[(LLF)] If $L_2$ satisfies Mutex and Live-Lock Freedom but not Starvation Freedom for two processes, then $L_n$ should satisfy Mutex and Live-Lock Freedom but not Starvation Freedom for $n$ processes.
        \item[(SF)] If $L_2$ satisfies Mutex and Starvation Freedom for two processes, then $L_n$ should satisfy Mutex and Starvation Freedom for $n$ processes.
    \end{itemize}
    You must design the reduction and prove that it satisfies the requirements (LLF) and (SF). You may assume that $n$ is fixed and known and that the processes are named $1, 2, \ldots, n$.
\end{problem}

\begin{solution}
    \bbni

    We can construct a lock $L_n$ that solves the mutual exclusion problem for $n$-processes by constructing a tournament bracket for the $n$ processes. 

    We construct a complete binary tree of height $h = \ceil{\log_2 n}$ with the first $n$ leaf nodes corresponding to processes $1, \ldots, n$ (the remaining leaves are unused). Each internal node acts as an independent $L_2$ lock, and its two roles $0$ and $1$ correspond to its left and right subtrees. Where if they `win' that node's $L_2$ lock, they move to the next node instead of entering CS. Also note that by `holds' we mean that a process has finished a node's Try section but hasn't yet started its Exit section.

\begin{tcolorbox}
     \begin{center}
    \begin{tikzcd}[column sep={0.8cm,between origins}, row sep={1.4cm,between origins}, arrows={dash}, every label/.append style={font=\scriptsize}]
        & & & & & & & |[lock, label=right:{\ensuremath{n_3\mbox{ (root)}}}]| L_2
            \arrow[dllll, "0"'] \arrow[drrrr, "1"] & & & & & & & \\
        & & & |[lock, label=left:{\ensuremath{n_2}}]| L_2
            \arrow[dll, "0"'] \arrow[drr, "1"] & & & & & & & & |[lock]| L_2
            \arrow[dll, "0"'] \arrow[drr, unused] & & & \\
        & |[lock]| L_2 \arrow[dl, "0"'] \arrow[dr, "1"] & & & & |[lock, label=right:{\ensuremath{n_1}}]| L_2
            \arrow[dl, "0"'] \arrow[dr, "1"] & & & & |[lock]| L_2
            \arrow[dl, "0"'] \arrow[dr, unused] & & & & |[lock, unused]| \phantom{L_2}
            \arrow[dl, unused] \arrow[dr, unused] & \\
        |[proc]| 1 & & |[proc]| 2 & & |[proc]| 3 & & |[proc]| 4 & & |[proc]| 5 & & |[proc, unused]| \phantom{0} & & |[proc, unused]| \phantom{0} & & |[proc, unused]| \phantom{0}
    \end{tikzcd}
    \end{center}
    {The tournament tree for $n = 5$, so $h = \ceil{\log_2 5} = 3$. Each circle is an independent $L_2$ lock, and the edge labels give the side (role) a process uses at that lock. Dashed nodes are unused.}
\end{tcolorbox}
   




    Let $n_1, \ldots, n_h$ be the path from a process $p$'s leaf nodes parent to the root node, and let $r_i(p) \in \{0, 1\}$ indicate which side of node $n_i$ contains $p$. Then for a process $p$ we define it's algorithm for the Try and Exit sections of the $L_n$ lock to be: 

\begin{lstlisting}[mathescape]
    Try$_n$($p$): 
        for $i$ in range$(1, h+1)$: 
            Try$_{{n_i}}$(r$_i$(p))

        enter CS
    
    Exit$_n$($p$): 
        for $i$ in range$(h, 0, -1)$: 
            Exit$_{{n_i}}$(r$_i$(p))
        return to Remainder
\end{lstlisting}
    Note that a process holds every `win' it acquires until it leaves the Critical Sectionm and that it releases its wins in reverse order, starting from the root and working back down to its leaf. We say a process holds a lock on a node if it has finished that node's Try section but has not yet finished its Exit Section. Every process in the CS of $L_n$ holds the root.  


    \textbf{Mutex} 

    We will show via induction on the height of the node that Mutex is satisfied. 
    
    Base Case (height $= 1$): For parents of leaves, each side only has one process, and hence by the Mutex of the $L_2$ of that node this is clear.

    Inductive Hypothesis (height $= h-1$): Suppose that for a higher node $v$, at most one process approaches it from each side (by Mutex of its child $c$). 
    
    Inductive Step: Let $v$ be a node of height $h$, and let $c$ be one of its child nodes. Any process on that side of $v$ already has a lock on $c$. Then, since the locks are released in reverse order (root to leaves), the process doesn't start Exit$_c$ until it has finished Exit$_v$, hence the process holds $c$ the whole time it is in $v$. By the inductive hypothesis, $c$'s $L_2$ lock is used as a valid two-process lock, so by Mutex of $c$'s $L_2$ at most one process can hold $c$. Thus at most one process can be using that side of $v$.

    Hence the root's $L_2$ lock is also valid and by its Mutex at most one process can hold the root, and thus $L_n$ satisfies Mutex.

    \bbni
    \textbf{(LLF)} 
    
    Suppose $L_2$ satisfies Live-Lock Freedom, so if some process is in an $L_2$ Try section and nobody holds that lock, then eventually some process wins it. Assume for contradiction that some process stays in Try$_n$ forever and that after some point no process ever enters the CS again. 

    Then once all previously started CS and Exit sections finish, no process ever releases a node again, since processes only release nodes after leaving the CS. Thus processes only move up the tree and each process can move up at most $h$ times. So after some point no one wins any more locks. Thus every process is either in the Remainder section or stuck forever in the Try section of some node. 

    Let $v$ be the highest node where a process $q$ is stuck in the Try section. Since no one wins the lock for $v$ again, by the Live Lock Freedom of $v$'s $L_2$ someone must be holding $v$ forever. Let that process be $s$. If $v$ is the root then $s$ eventually enters and releases the CS, a contradiction. Otherwise, $s$ holds $v$ but never releases it, so $s$ is not in the CS and must be stuck at some node above $v$. This contradicts our choice of $v$ as the highest node where a process is stuck in the Try section.

    Thus some process eventually enters the CS, so $L_n$ satisfies Live-Lock Freedom. 
    To show that it is not starvation-free when $L_2$ is not, let only two processes try, one from each side of the root. They can get their lower locks since nobody else is trying for those locks. At the root, let them take steps in the same order as an $L_2$ execution where one process starves. The extra steps for the lower locks all finish, so the execution is still fair. Thus same process waits forever at the root and never enters the CS, and $L_n$ is not starvation-free. 


    \bbni
    \textbf{(SF)} 
    
    
    Suppose $L_2$ satisfies Starvation Freedom, so any process in an $L_2$ Try section eventually wins as long as whoever holds that lock eventually releases it. 

    Let $p$ be any process in Try$_n$ with path $n_1, \ldots, n_h$. We will show that if $p$ is in the Try section of $n_i$, then $p$ eventually wins $n_i$, and eventually enters the CS via induction on $i$ that going down from the root $(i = h)$ to $i=1$. 

    Base Case ($i=root$): At the root, whatever process is using the CS eventually is released. Note that Exit$_n$ releases the root first. Then by the starvation freedom of the $L_2$ of the root $p$ eventually wins it and enters the CS. 
    
    Inductive Hypothesis: Suppose that the claim holds for every node $n_j$ for $j>i$. 

    Inductive Step: Let $q$ be a process holding $n_i$. Since $q$ is above $n_i$, it is in one of three places: 
    \begin{enumerate}
        \item in the Try section of some $n_j$ with $j > i$. Then by IH $q$ eventually enters the CS. 
        \item in the CS, by starvation freedom of the $L_2$ of the root $q$ eventually leaves it. 
        \item in Exit$_n$. 
    \end{enumerate}

    In every case $q$ eventually reaches Exit$_n$, which releases its locks from the root down to $n_i$, so whoever holds $n_i$ always releases it, and by Starvation Freedom of $n_i$'s $L_2$ lock, $p$ eventually wins $n_i$. Then $p$ moves on to the Try section of $n_{i+1}$, and by the inductive hypothesis $p$ eventually enters the CS.

    Thus any process that starts Try$_n$ eventually wins every node on its path and enters the CS. Hence $L_n$ satisfies Starvation Freedom.
    

    
\end{solution}


\end{document}